%-------------------------------------------------------------------------------
% Experience — STEM / Research Tools variant (builder-forward)
% Also used by the Research variant (cv-research.tex).
% Hello Soda, Consultant and USC entries are shared: edit them in
% cv-sections/shared/.
%-------------------------------------------------------------------------------
\cvsection{Experience}

\begin{cventries}

\cventry
{CEO of Impact}
{Quansight PBC}
{Manchester, UK}
{May 2025 -- Present}
{
\begin{cvitems}
  \item {Maintain, design, write and review tooling used by millions of developers and researchers every day; currently improving conda's environment resolution and package verification so environments are reproducible and install faster.}
  \item {Lead partnerships that fund maintainer time on core scientific computing work: free-threaded Python with Meta, documentation, tooling and demos across JAX and OpenXLA with Google, and ecosystem-wide packaging improvements with NVIDIA, Red Hat and Meta.}
  \item {Set technical priorities across our open source portfolio (conda, conda-forge, Jupyter, Nebari), weighing what scientists, maintainers and clients need most. Brought in new partners including Red Hat, DeepMind and the Gates Foundation.}
  \item {Answer to the board for both sides of Quansight's mission as a Public Benefit Corporation: a sustainable business and a healthier open source ecosystem.}
\end{cvitems}
}

\cventry
{Director}
{Quansight Labs}
{Manchester, UK}
{January 2021 -- May 2025}
{
\begin{cvitems}
  \item {Took over Nebari, an open source platform for deploying AI and ML infrastructure, and moved it from an internal tool to a community project with its own governance. As product owner for around two years, set the roadmap, helped design the architecture, and wrote, reviewed and tested code for GPU provisioning, model serving and reproducible environments, used by client teams and research organisations including NOAA and NASA.}
  \item {Architected and led a computer-vision ML platform, built from scratch, that detects brands and logos across sports, esports and live-event video and measures their on-screen visibility. The client's own team used it to train and onboard new brand models continuously; it scaled to 10,000+ brands and 8+ hours of video a day, at 40\% lower cost than the third-party product it replaced.}
  \item {Led a 50+ person team of software engineers, data scientists, ML engineers, UX designers and developer-experience engineers, setting technical direction across concurrent client and open source engagements.}
  \item {Grew Quansight Labs' open source work 200\% over the role, measured across staff hours, pull requests, community work such as calls and mentoring, and the number of projects we supported; aligned it with our consultancy work, growing revenue 25\%.}
  \item {Designed and tested ways to keep open source sustainable (consortia, open core, collaborative funding, community-driven models) and advised other organisations on open source strategy and governance.}
\end{cvitems}
}

\cventry
{Sr. Developer Advocate \& Tech Lead}
{Microsoft}
{Manchester, UK}
{March 2019 -- December 2020}
{
\begin{cvitems}
  \item {Brought feedback from the scientific and numerical Python communities into VS Code, working with the Jupyter notebook and Pylance teams on better plotting in notebooks (Matplotlib, Bokeh and other interactive plots), Python environment discovery, the Gather extension, and many smaller usability fixes.}
  \item {Contributed features and performance improvements to VS Code, Azure Functions and the Azure SDK, and worked with Azure leadership on competitive analysis that shaped the Azure Functions CLI and SDK experience.}
  \item {Worked with Microsoft Research on user experience and optimisation improvements for InnerEye, CNTK and Azure ML.}
  \item {Co-led a cloud platform for the UK's national weather service, processing millions of observations a day for 10+ research teams.}
  \item {Partnered with the Alan Turing Institute to build a secure, scalable data platform used by hundreds of data scientists across the UK.}
  \item {Co-developed Microsoft's Python advocacy strategy for scientific computing, data science and education, and delivered 120+ talks, workshops and keynotes.}
\end{cvitems}
}

%-------------------------------------------------------------------------------
% Hello Soda — shared experience entry, \input inside each variant's cventries
%-------------------------------------------------------------------------------
\cventry
{Sr. Machine Learning Engineer}
{Hello Soda}
{Manchester, UK}
{July 2018 -- February 2019}
{
\begin{cvitems}
  \item {Architected a platform running 60+ ML and NLP models for customers worldwide, improving model uptime and performance by 15\%.}
  \item {Introduced MLOps practices and automation that cut model deployment from 14 days to 2 and improved system reliability by 40\%.}
  \item {Put data governance and deployment standards in place that cut operational costs 25\%, improved model accuracy 20\% and reduced production incidents.}
\end{cvitems}
}

%-------------------------------------------------------------------------------
% Independent consultant — shared experience entry
%-------------------------------------------------------------------------------
\cventry
{Independent Consultant}
{Freelance}
{Manchester, UK}
{November 2016 -- Present}
{
\begin{cvitems}
  \item {Advise start-ups, researchers and organisations on reproducible data workflows, data pipelines from ETL to deployment, and data and open source governance, reducing computation time on data-heavy workflows in most engagements.}
\end{cvitems}
}


\cventry
{Research Software Engineer, ML Team}
{The University of Sheffield}
{Sheffield, UK}
{November 2016 -- June 2018}
{
\begin{cvitems}
  \item {Built analysis pipelines in Python and R for NLP, bioinformatics, materials science and geospatial research, from data collection and wrangling to models and results.}
  \item {Built dashboards and visualisations with d3.js, Shiny, Bokeh and Plotly so researchers could explore and share their findings.}
  \item {Core contributor to OpenDreamKit (EU-funded) and Project Jupyter, building tools for reproducible research and literate programming.}
  \item {Mentored colleagues and students in reproducible workflows, software engineering and HPC, and served on the UK RSE and NumFOCUS diversity and inclusion committees.}
\end{cvitems}
}

%-------------------------------------------------------------------------------
% University of Southern California — shared experience entry
%-------------------------------------------------------------------------------
\cventry
{Research Engineer, Medical Devices}
{University of Southern California}
{California, USA}
{2011 -- 2012}
{}


\end{cventries}
